\documentclass[11pt,a4paper]{article}

\usepackage[utf8]{inputenc}
\usepackage[T1]{fontenc}
\usepackage[english]{babel}
\usepackage{amsmath,amssymb,amsthm}
\usepackage{geometry}
\usepackage{hyperref}
\usepackage{listings}
\usepackage{xcolor}
\usepackage{booktabs}
\usepackage{algorithm}
\usepackage{algpseudocode}
\usepackage{graphicx}

\geometry{margin=2.5cm}
\hypersetup{colorlinks=true,linkcolor=blue,citecolor=blue,urlcolor=blue}

\title{A Hybrid Knowledge-Graph RAG Pipeline for Legal PDF Corpora:\\
Integrating ODKE+ Extraction with KG-GPT and ARK-V1 Reasoning}
\author{Samuel Bagin}
\date{\today}

\begin{document}
\maketitle

\begin{abstract}
We present the design and implementation of \texttt{PDFGraphRAG}, a Python system
that transforms Slovak-language legal PDFs into a Neo4j knowledge graph and
answers natural-language questions over it. The ingestion stage follows the
\emph{Open-Domain Detection $\rightarrow$ Schema Refinement $\rightarrow$
Schema-Driven Extraction} (ODD/SR/SDE) paradigm popularised by
ODKE+~\cite{odkeplus2025}, augmented with a document-layout-aware table
extractor based on DocLayout-YOLO. The query stage follows the three-stage
KG-GPT pipeline~\cite{kggpt2023}---sentence segmentation, graph retrieval and
inference---but replaces the original LLM ranker of Stage~2 with the
ARK-V1~\cite{arkv12025} state-machine retriever, which grounds every relation
selection in the live graph schema. This document formalises the pipeline,
describes each module of \texttt{code/pdf\_graphrag.py}, and relates the design
choices to the three source papers using mathematical notation.
\end{abstract}

\section{Introduction}

Retrieval-Augmented Generation (RAG) over unstructured text is now standard,
yet vector similarity alone collapses the rich relational structure implicit
in legal documents: paragraphs, sections, obligations and temporal
constraints. Knowledge-graph RAG (KG-RAG) re-introduces that structure by
grounding an LLM's answers in a typed graph
$G=(V,E,\mathcal{T}_V,\mathcal{T}_E)$ with node labels
$\mathcal{T}_V$ and relation labels $\mathcal{T}_E$. Our system combines three
recent lines of work:

\begin{enumerate}
    \item \textbf{ODKE+}~\cite{odkeplus2025}, which advocates an ontology-guided
    extraction loop with a distinct \emph{grounder} that filters
    hallucinations, reporting a 35\% reduction in unsupported facts.
    \item \textbf{KG-GPT}~\cite{kggpt2023}, which decomposes KG question
    answering into (\textsc{Segment}, \textsc{Retrieve}, \textsc{Infer}).
    \item \textbf{ARK-V1}~\cite{arkv12025}, a state-machine agent that iterates
    (\textsc{relation-enum}, \textsc{relation-select}, \textsc{triple-enum},
    \textsc{reason}) steps up to a bounded depth $K_{\max}$.
\end{enumerate}

Section~\ref{sec:ingestion} describes the ingestion phase, Section~\ref{sec:query}
the query phase, Section~\ref{sec:impl} discusses implementation details, and
Section~\ref{sec:discussion} places the design against the cited literature.

\section{Ingestion: From PDF to Typed Graph}
\label{sec:ingestion}

Let a document $D$ be a sequence of pages $(p_1,\dots,p_n)$. The ingestion
pipeline produces two artefacts: a graph $G_D\subseteq G$ and a set of vector
indices $\{\mathbf{V}_N,\mathbf{V}_C,\mathbf{V}_R\}$ over nodes, chunks and
relation-type labels respectively. The top-level driver is
\texttt{PDFGraphRAG.process} in \texttt{code/pdf\_graphrag.py}.

\subsection{Layout-aware table extraction}

Before textual processing, \texttt{detect\_tables} rasterises the PDF at
200~DPI and runs a DocLayout-YOLO v10 model (checkpoint
\texttt{juliozhao/DocLayout-YOLO-DocStructBench}) with a confidence threshold
$\tau=0.3$. For each page $p_i$ the detector returns a set of bounding boxes
$B_i=\{(x_1,y_1,x_2,y_2,c)\}$ where $c\in\mathcal{C}_{\text{target}}=\{\text{table},
\text{picture}, \text{figure}, \text{isolate\_formula}, \text{formula\_caption}\}$.
Consecutive table detections (page distance $\leq 1$) are grouped by
\texttt{group\_table\_detections} into multi-page tables. Each group is sent
to \texttt{transform\_table\_to\_html}, which prompts a vision-enabled
Gemini 2.5 Pro model with the concatenated page crops and a strict
structured-output schema
\texttt{TableHTMLResponse} $=(\text{html}, \text{page\_range}, \text{row\_count},
\text{column\_count})$. The resulting HTML is then transformed into a
hierarchical subgraph by \texttt{transform\_html\_to\_graph\_document}, which
instructs the LLM to emit a rooted tree whose nodes satisfy the invariants
\[
\forall v\in V_{\text{table}}:\ \text{id}(v)\text{ unique}\wedge
\text{name}\in\text{props}(v)\wedge \text{ASCII-normalised}(v).
\]
Pages that are strictly interior to a multi-page table are removed from the
text stream to avoid duplicate extraction, while the first and last pages are
retained because they may contain surrounding prose.

\subsection{Open-Domain Detection (ODD)}

ODD mirrors the ``schema discovery'' phase of ODKE+~\cite{odkeplus2025}, in
which a language model proposes candidate predicates before ontology pruning.
The text is split with a \texttt{RecursiveCharacterTextSplitter} of chunk size
$s=1200$ and overlap $o=200$. For each chunk $c_i$ an asynchronous agent
(\texttt{open\_domain\_detection}) returns a local schema
\[
S_i=(\mathcal{T}^{(i)}_V,\mathcal{T}^{(i)}_E),
\]
by calling \texttt{create\_agent} with a \texttt{ProviderStrategy} that forces
adherence to \texttt{response\_schema\_for\_odd}. Concurrency is controlled by
an \texttt{asyncio.Semaphore($M$)} (default $M=5$) and a cooperative
pause/retry event that back-offs the whole batch for 60\,s on provider errors,
matching the robustness pattern recommended by ODKE+ for production
throughput~\cite{odkeplus2025}. The aggregated schema is
\[
S_{\text{ODD}}=\Big(\bigcup_i \mathcal{T}^{(i)}_V,\ \bigcup_i \mathcal{T}^{(i)}_E\Big).
\]

\subsection{Schema Refinement (SR)}

The aggregated schema is still noisy: diacritic variants (e.g., \emph{Dan}
vs.\ \emph{Daň}), near-synonyms and ontology mismatches are common.
\texttt{schema\_refinement} invokes Gemini 2.5 Pro with a structured-output
schema \texttt{SchemaRefinementResponse} that returns canonical label sets
together with a merge log $\mathcal{M}:S_{\text{ODD}}\to S^{\ast}$. The
normalisation is conservative---the prompt explicitly forbids merging
ontologically distinct kinds (person vs.\ document) and enforces
\textsc{UPPER\_SNAKE\_CASE} for relations. Following
ODKE+~\cite{odkeplus2025}, an \emph{existing ontology} $S_0$ (the current
Neo4j schema) is passed in so that the refined schema $S^\ast$ remains
compatible with facts already stored:
\[
S^{\ast}=\Phi(S_{\text{ODD}},\,S_0),\qquad
\Phi=\text{LLM-refiner}.
\]

\subsection{Schema-Driven Extraction (SDE)}

With $S^\ast$ fixed, the text is resplit with $s=1024$, $o=128$ and processed
by \texttt{async\_schema\_driven\_extraction}. Each chunk $c_i$ yields a
\texttt{GraphDocument}
$G_i=(V_i,E_i,\text{src}_i)$ via \texttt{\_convert\_to\_graph\_document}.
This converter enforces four invariants that together implement an
in-pipeline \emph{grounder} analogous to the one advocated by
ODKE+~\cite{odkeplus2025}:
\begin{enumerate}
  \item \textbf{ID validation}: nodes without a non-empty \texttt{id} are
  dropped (line~441 of \texttt{pdf\_graphrag.py}).
  \item \textbf{Type fallback}: \texttt{format\_node\_type} defaults empty
  labels to \texttt{DEFAULT\_NODE\_TYPE} = \texttt{Entity}.
  \item \textbf{Property keying}: \texttt{format\_property\_key} converts keys
  to camelCase; \texttt{format\_relationship\_type} normalises relations to
  uppercase-underscore.
  \item \textbf{Relationship grounding}: a relation is emitted only if both
  endpoints resolve to a node of the same chunk (case-insensitive match), and
  a \texttt{Chunk} node is linked to every extracted node via
  \texttt{IN\_CHUNK}.
\end{enumerate}

An optional \texttt{\_filter\_by\_strict\_mode} post-filters the graph to the
allowed label sets, preserving \texttt{Chunk} nodes and \texttt{IN\_CHUNK}
edges by construction.

\subsection{Persistence and vector stores}

The graph documents---both from text (SDE) and from tables---are inserted via
\texttt{Neo4jGraph.add\_graph\_documents} with \texttt{baseEntityLabel=True}.
Two Cypher statements then strip the auxiliary label from \texttt{Chunk} and
\texttt{Document} nodes so that the entity vector index excludes them. Three
Neo4jVector indices are built with OpenAI \texttt{text-embedding-3-large}:
\[
\mathbf{V}_N\text{ over nodes},\quad
\mathbf{V}_C\text{ over chunks},\quad
\mathbf{V}_R\text{ over relation-type labels}.
\]
The relation-type index is populated by querying
\texttt{CALL db.relationshipTypes()} and embedding each label as its own
pseudo-document, enabling later relation-matching at query time.

\section{Query: Three-Stage KG-GPT with ARK-V1 Retrieval}
\label{sec:query}

Given a question $q$, \texttt{PDFGraphRAG.query} implements the canonical
KG-GPT pipeline~\cite{kggpt2023}
\[
q \xrightarrow{\text{Stage 1}} \{s_j\}_{j=1}^m
  \xrightarrow{\text{Stage 2}} \mathcal{E}
  \xrightarrow{\text{Stage 3}} a,
\]
where $s_j$ are sub-sentences, $\mathcal{E}$ is a set of evidence triples and
$a$ is a natural-language answer.

\subsection{Stage 1: Sentence segmentation}

\texttt{segment\_question} prompts \texttt{openai\_client} with
\texttt{system\_prompt\_for\_segmentation}. The agent emits a
\texttt{response\_schema\_for\_segmentation} object parsed into
\texttt{SubSentence} records:
\[
s_j=(\text{text}_j,\ \text{entities}_j),\quad |\text{entities}_j|\leq 2.
\]
The constraint of at most two anchor entities per sub-sentence enforces the
KG-GPT assumption~\cite{kggpt2023} that every sub-sentence corresponds to a
\emph{single implied triple} $(h,r,t)$. If segmentation fails, the pipeline
falls back to a one-element list $\{(q,\emptyset)\}$ to preserve liveness.

\subsection{Stage 2: ARK-V1 state-machine retrieval}

Whereas KG-GPT's original Stage~2 scores a static candidate set of relations,
our implementation uses the ARK-V1~\cite{arkv12025} state-machine, realised by
\texttt{ark\_v1\_retrieve}. The state is a 4-tuple
\[
\sigma_k=(a_k,\ \text{summary}_k,\ R^{(k)},\ T^{(k)}),
\]
where $a_k$ is the current anchor, $\text{summary}_k$ is a running NL
summary of prior steps, $R^{(k)}\subseteq\mathcal{T}_E$ are the relations
actually present on $a_k$ in the graph, and $T^{(k)}$ is the candidate triple
set. A step proceeds as follows:

\begin{enumerate}
  \item \textbf{Anchor resolution} (\texttt{\_resolve\_anchor}): an entity
  mention is lifted to an actual graph node by case-insensitive exact match,
  with a substring fallback.
  \item \textbf{Relation enumeration} (\texttt{\_retrieve\_relations}):
  $R^{(k)}=\{\,\text{type}(r)\mid \exists u.\,(a_k,r,u)\in E\text{ or }(u,r,a_k)\in E\,\}$.
  \item \textbf{Relation selection} (\texttt{\_select\_relation}):
  \texttt{gemini\_client\_flash} with a \texttt{RelationChoice} schema picks
  $r_k\in R^{(k)}\cup\{\bot\}$. Up to $C_{\max}=2$ retries re-prompt the LLM
  if it hallucinates a relation outside $R^{(k)}$; this is exactly the
  \emph{grounded relation selection} invariant of ARK-V1~\cite{arkv12025}.
  \item \textbf{Triple enumeration} (\texttt{\_retrieve\_triples}):
  $T^{(k)}=\{(h,r_k,t)\in E\}$ limited to 25 rows per direction.
  \item \textbf{Reasoning} (\texttt{\_reason\_step}): an LLM returns a
  \texttt{ReasoningStep} object
  \[
  \rho_k=(I_k\subseteq\{0,\dots,|T^{(k)}|-1\},\ \text{impl}_k,\
  \text{cont}_k\in\{0,1\},\ a_{k+1}\in\{t\,:\,(h,r_k,t)\in T^{(k)}\}\cup\{\bot\}).
  \]
  Selected triples are appended to $\mathcal{E}$ and their endpoints to the
  node list used later for chunk retrieval.
\end{enumerate}

The loop terminates when (i) $R^{(k)}=\emptyset$, (ii) the LLM emits $\bot$,
(iii) \texttt{cont}$_k=0$, (iv) the proposed $a_{k+1}$ is not a tail of a
selected triple, or (v) the depth cap $k=K_{\max}=3$ is reached. The final
evidence set is de-duplicated while preserving insertion order:
\[
\mathcal{E}_{\text{final}}=\text{dedup}\Big(\bigcup_{j}\bigcup_{k}\{t:t\in T^{(k)},\ \text{idx}(t)\in I_k\}\Big).
\]

Two design details deserve emphasis. First, the prompt for
\texttt{\_select\_relation} is reconstructed from scratch every step
(system~+~sub-sentence~+~summary~+~current $R^{(k)}$), so the token budget is
$\mathcal{O}(1)$ in $k$ rather than growing with the reasoning trace, exactly
as recommended by Klein and Ohnemus~\cite{arkv12025}. Second, every Cypher
query passes through \texttt{is\_read\_only\_cypher}, which rejects write
keywords and the \texttt{CALL} procedure family via pre-compiled regular
expressions, giving the agent a capability-restricted tool surface.

\subsection{Chunk retrieval via \texttt{IN\_CHUNK}}

After Stage~2, \texttt{get\_chunks\_from\_nodes} issues
\begin{verbatim}
MATCH (n)-[:IN_CHUNK]->(c:Chunk)
WHERE n.id IN $node_ids
RETURN DISTINCT c.text, c.page, n.id
ORDER BY c.page
\end{verbatim}
to recover the original text passages from which each retrieved entity was
extracted. This corresponds to the KG-GPT~\cite{kggpt2023} observation that
symbolic triples alone can be too terse to disambiguate numeric or temporal
constraints, and that attaching source text to the prompt substantially
improves final-answer faithfulness.

\subsection{Stage 3: Inference}

\texttt{answer} formats the evidence as a JSON list
$\mathcal{E}_{\text{final}}=\{[h,r,t]\}$ and the chunks as a
\texttt{\textbackslash n---\textbackslash n}-delimited string, then invokes
\texttt{gemini\_client\_flash} with a structured schema
\texttt{InferenceAnswer} so that the model's output is forced into a single
\texttt{answer} field. The prompt explicitly instructs the model to treat the
triples as \emph{primary evidence} and the chunks as \emph{supplementary
context}, which is the precedence rule proposed in KG-GPT's Stage-3
ablations~\cite{kggpt2023}. Schematically,
\[
a=\text{LLM}\big(q,\ \mathcal{E}_{\text{final}},\ \mathcal{C}_{\text{final}}\big),
\]
where $\mathcal{C}_{\text{final}}$ is the chunk context.

\section{Implementation Notes}
\label{sec:impl}

\subsection{LLM router}

The constructor of \texttt{PDFGraphRAG} instantiates four clients: an OpenAI
model for agentic tasks (segmentation, ODD, SDE, search-agent), two Gemini
2.5 Pro variants (baseline and high-thinking) for heavier reasoning, a
Gemini 2.5 Flash for lightweight schema-constrained calls
(relation selection, reasoning step, final answer), and an Anthropic client
kept for future use. Cost/latency arguments thus drive the routing: Flash for
per-step agent calls inside the ARK-V1 loop (latency-sensitive, small
context), Pro for one-shot refinement and table-to-graph extraction (higher
accuracy at larger context).

\subsection{Serialisation}

\texttt{serialize\_for\_json} recognises Neo4j \texttt{Node},
\texttt{Relationship} and \texttt{Path} objects by duck-typing
(\texttt{hasattr(obj,'labels')} etc.) and recursively converts them to
JSON-safe dictionaries. This is required because the
\texttt{search\_database} tool inside \texttt{search\_agent\_retrieve}
returns native driver objects that the Langchain tool contract must
stringify.

\subsection{Idempotent graph writes}

\texttt{add\_graph\_docs\_without\_apoc} issues \texttt{MERGE} on both nodes
and relationships, so re-ingesting the same PDF does not duplicate edges. The
main \texttt{process} method prefers \texttt{graph.add\_graph\_documents} with
\texttt{baseEntityLabel=True}, which depends on APOC and is therefore a
deployment constraint for the target Neo4j instance.

\section{Discussion}
\label{sec:discussion}

\paragraph{Relation to ODKE+.}
ODKE+~\cite{odkeplus2025} targets Wikipedia-scale fact freshness with a
five-stage pipeline (Initiator, Retriever, Extractor, Grounder, Corroborator)
and reports 98.8\% precision after corroboration. Our ingestion adopts the
\emph{Extractor} and \emph{Grounder} abstractions: the ODD/SR/SDE triple is a
specialisation of ODKE+'s ontology-guided extractor for the case where the
ontology is only \emph{partially} known (we bootstrap it from the corpus),
while \texttt{\_convert\_to\_graph\_document} and
\texttt{\_filter\_by\_strict\_mode} together play the role of the grounder.
Unlike ODKE+ we do not yet implement a \emph{corroborator} stage---a natural
extension would be cross-chunk agreement scoring before persistence.

\paragraph{Relation to KG-GPT.}
Kim et al.~\cite{kggpt2023} split KGQA into three LLM-driven stages and show
that factoring the problem this way outperforms monolithic prompting on
MetaQA and FactKG. We reproduce their factorisation verbatim at the level
of the \texttt{query} method, and we additionally import their observation
that attaching raw text context to Stage~3 improves numeric/temporal
reasoning. The novel element is the substitution of Stage~2: KG-GPT's
Lexical-Matching-and-LLM-Ranking (LMLR) retriever is replaced by ARK-V1.

\paragraph{Relation to ARK-V1.}
Klein and Ohnemus~\cite{arkv12025} formalise the retriever as a Mealy machine
whose transitions are gated by LLM calls constrained to the empirically
observed relation set $R^{(k)}$. Our implementation follows the paper
closely: bounded depth, bounded retries on invalid LLM outputs, and a
constant-size prompt that elides older intermediate steps behind an NL
summary. Deviations are minor: we bound anchor expansion to the first two
entity mentions (\texttt{(entities or [])[:2]}) and the triple limit to
25 per direction---both pragmatic caps that keep Gemini Flash calls below
the free-tier token ceiling.

\paragraph{Limitations.}
(i)~The grounder uses surface-form validation only; a semantic grounder
(entailment classifier) as in ODKE+~\cite{odkeplus2025} would further
reduce hallucinations. (ii)~ARK-V1 with $K_{\max}=3$ under-covers genuine
multi-hop questions; adaptive depth selection conditioned on
sub-sentence complexity would be an upgrade. (iii)~Slovak diacritics are
stripped at extraction time; a post-hoc re-diacritisation step would
improve user-facing rendering of the final answer.

\section{Conclusion}

\texttt{PDFGraphRAG} is, to our knowledge, the first open implementation
that composes ODKE+-style hybrid extraction~\cite{odkeplus2025} with the
KG-GPT three-stage query pipeline~\cite{kggpt2023} and the ARK-V1
state-machine retriever~\cite{arkv12025} in a single Python module,
targeted at Slovak-language legal corpora. The ingestion side
contributes a layout-aware table extractor that preserves multi-page
table hierarchies as typed subgraphs, and the query side contributes an
\texttt{IN\_CHUNK}-based dual retrieval that feeds both symbolic triples
and source prose to the final LLM. Empirical evaluation on a held-out
Slovak tax-law benchmark is left to future work.

\begin{thebibliography}{9}

\bibitem{odkeplus2025}
Khorshidi, S.\ et al.\ (Apple Inc.).
\newblock \emph{ODKE+: Ontology-Driven Knowledge Extraction at Web Scale}.
\newblock arXiv:2509.04696, 2025.
\newblock \url{https://arxiv.org/html/2509.04696v1}.

\bibitem{kggpt2023}
Kim, J., Kwon, Y., Jo, Y., and Choi, E.
\newblock \emph{KG-GPT: A General Framework for Reasoning on Knowledge Graphs
Using Large Language Models}.
\newblock Findings of EMNLP 2023.
\newblock \url{https://aclanthology.org/2023.findings-emnlp.631.pdf}.

\bibitem{arkv12025}
Klein, J.\ F.\ and Ohnemus, J.
\newblock \emph{ARK-V1: A State-Machine LLM Agent for Knowledge-Graph
Question Answering}.
\newblock arXiv:2509.18063, 2025.
\newblock \url{https://arxiv.org/pdf/2509.18063}.

\end{thebibliography}

\end{document}
